%%%% Начало оформления заголовка - оставить без изменений !!! %%%%
\input{chapters/task_settings}	% настройки - начало


{\centering%

	\Ministry\\
	\SPbPU\\
	\institute{}
\par}
}\intervalS{}% завершает input

				\noindent
				\begin{minipage}{\linewidth}
				\vspace{\mfloatsep} % интервал
				\begin{tabularx}{\linewidth}{Xl}
					&УТВЕРЖДАЮ\\
					&\HeadTitle{}\\
					&\underline{\hspace*{0.1\textheight}} \Head{}\\
					&<<\underline{\hspace*{0.05\textheight}}>> \underline{\hspace*{0.1\textheight}} \approveYear~г.\\
				\end{tabularx}
				\vspace{\mfloatsep} % интервал
				\end{minipage}

\intervalS{\centering\bfseries%

				ЗАДАНИЕ\\
				на выполнение %с 2020 года
				%по выполнению % до 2020 года
				выпускной квалификационной работы


\intervalS\normalfont%

				студенту \uline{\AuthorFullDat{}, гр.~\group}


\par}\intervalS%
%%%%
%%%% Конец оформления заголовка  %%%%



\begin{enumerate}[1.]
	\item Тема работы: {\expandafter{}\ulined{} <<\thesisTitle>>.}
	%\item Тема работы (на английском языке): \uline{\thesisTitleEn.} % вероятно после 2021 года
	\item Срок сдачи студентом законченной работы: \uline{\thesisDeadline.}
	\item Исходные данные по работе: \uline{спецификации OAuth~2.0~\cite{RFC6749}, OpenID Connect~\cite{OpenIDConnect}, TLS~1.3~\cite{RFC8446}; научные публикации по~теме~\cite{Almeida2022, Barabanov2021}.}
	\printbibliographyTask{} % печать списка источников
	\item Содержание работы (перечень подлежащих разработке вопросов):
	\begin{enumerate}[label=\theenumi\arabic*.]
		\item Обзор и~анализ подходов к~обеспечению безопасности в~распределённых микросервисных системах.
		\item Проектирование архитектуры AAA-подсистемы.
		\item Реализация компонентов пользовательской аутентификации и~авторизации.
		\item Реализация компонентов межсервисной аутентификации и~авторизации.
		\item Реализация компонентов аудита и~мониторинга.
		\item Апробация разработанной AAA-подсистемы.
	\end{enumerate}
	\item Перечень графического материала (с указанием обязательных чертежей):
	\begin{enumerate}[label=\theenumi\arabic*.]
		\item Архитектурная схема AAA-подсистемы.
		\item Диаграмма взаимодействия компонентов при аутентификации и~авторизации.
		\item Схема развёртывания подсистемы в~среде Kubernetes.
	\end{enumerate}
		\item Перечень используемых информационных технологий, в том числе программное обеспечение, облачные сервисы, базы данных и прочие сквозные цифровые технологии (при наличии):
		\begin{enumerate}[label=\theenumi\arabic*.]
			\item Go~--- язык программирования.
			\item Docker, Kubernetes~--- контейнеризация и~оркестрация.
			\item TimeWeb Cloud~--- облачная платформа.
		\end{enumerate}
		\item Консультанты по работе:
		\begin{enumerate}[label=\theenumi\arabic*.]
		\item  \uline{\emakefirstuc{\ConsultantExtraDegree}, \ConsultantExtra.} % закомментировать при необходимости, идёт первый по порядку.
		\item \uline{\emakefirstuc{\ConsultantNormDegree}, \ConsultantNorm{} (нормоконтроль).} %	Обязателен для всех студентов
	\end{enumerate}
		\item Дата выдачи задания: \uline{\thesisStartDate.}
\end{enumerate}

\intervalS%можно удалить пробел

Руководитель ВКР \uline{\hspace*{0.1\textheight} \Supervisor}


\intervalS%можно удалить пробел

Консультант  \uline{\hspace*{0.1\textheight}\ConsultantExtra}


\intervalS%можно удалить пробел

%Консультант по нормоконтролю \uline{\hspace*{0.1\textheight} \ConsultantNorm}%ПОКА НЕ ТРЕБУЕТСЯ, Т.К. ОН У ВСЕХ ПО УМОЛЧАНИЮ

Задание принял к исполнению \uline{\thesisStartDate}

\intervalS%можно удалить пробел

Студент \uline{\hspace*{0.1\textheight}  \Author}



\input{chapters/task_settings_restore}	% настройки - конец
